% Quantum Mirror Theory - Supplementary Material
% 
% This file contains additional mathematical derivations and proofs
% that support the main paper.

\documentclass[12pt,a4paper]{article}

\usepackage[utf8]{inputenc}
\usepackage[T1]{fontenc}
\usepackage{amsmath,amssymb,amsfonts}
\usepackage{physics}
\usepackage{geometry}
\usepackage{hyperref}
\usepackage{booktabs}

\geometry{margin=1in}

\title{Supplementary Material:\\Quantum Mirror Theory}
\author{Babatope Jesse Afolabi\\
\textit{cr8OS Foundation DBA WPWakanda, LLC}}
\date{January 2026}

\begin{document}

\maketitle

\section{Mathematical Derivations}

\subsection{Proof: Mirror Operator is an Involution}

\textbf{Claim:} The Mirror Operator $M$ satisfying $M^2 = I$ is an involution.

\textbf{Proof:}

An involution is a function that is its own inverse. For $M$:
\begin{align}
M^2 &= I \\
M \cdot M &= I \\
M &= M^{-1}
\end{align}

Therefore, $M$ is self-inverse. Applying $M$ twice returns identity:
\begin{equation}
M(M\ket{\Psi}) = M\ket{\Psi'} = \ket{\Psi}
\end{equation}

You are your reflection's reflection. $\square$

\subsection{Eigenvalues of the Mirror Operator}

\textbf{Claim:} The eigenvalues of $M$ are $\pm 1$.

\textbf{Proof:}

Let $\ket{\lambda}$ be an eigenstate of $M$ with eigenvalue $\lambda$:
\begin{equation}
M\ket{\lambda} = \lambda\ket{\lambda}
\end{equation}

Applying $M$ again:
\begin{equation}
M^2\ket{\lambda} = M(\lambda\ket{\lambda}) = \lambda M\ket{\lambda} = \lambda^2\ket{\lambda}
\end{equation}

But $M^2 = I$, so:
\begin{equation}
I\ket{\lambda} = \ket{\lambda} = \lambda^2\ket{\lambda}
\end{equation}

Therefore $\lambda^2 = 1$, giving $\lambda = \pm 1$. $\square$

\subsection{Mirror Constant Range}

\textbf{Claim:} $0 \leq \mathbb{M} \leq 1$ for normalized states.

\textbf{Proof:}

The Mirror Constant is defined as:
\begin{equation}
\mathbb{M} = |\braket{\Psi|M|\Psi'}|
\end{equation}

By the Cauchy-Schwarz inequality:
\begin{equation}
|\braket{\Psi|M|\Psi'}| \leq \|\ket{\Psi}\| \cdot \|M\ket{\Psi'}\|
\end{equation}

For normalized states $\|\ket{\Psi}\| = \|\ket{\Psi'}\| = 1$, and since $M$ is unitary ($MM^\dagger = I$):
\begin{equation}
\|M\ket{\Psi'}\| = \|\ket{\Psi'}\| = 1
\end{equation}

Therefore:
\begin{equation}
\mathbb{M} \leq 1
\end{equation}

And by definition of absolute value:
\begin{equation}
\mathbb{M} \geq 0
\end{equation}

Thus $0 \leq \mathbb{M} \leq 1$. $\square$

\section{Shannon-Mirror Correspondence}

\subsection{Derivation}

Shannon entropy for a discrete random variable:
\begin{equation}
H = -\sum_{i} p_i \log_2 p_i
\end{equation}

For highly structured data (repeating patterns), the probability distribution is concentrated:
\begin{equation}
p_j \approx 1, \quad p_{i \neq j} \approx 0 \implies H \approx 0
\end{equation}

For random data (uniform distribution):
\begin{equation}
p_i = \frac{1}{n} \implies H = \log_2 n = H_{\max}
\end{equation}

The Mirror Constant maps inversely:
\begin{equation}
\mathbb{M} = 1 - \frac{H}{H_{\max}}
\end{equation}

Therefore:
\begin{itemize}
    \item $H = 0 \implies \mathbb{M} = 1$ (perfect coherence)
    \item $H = H_{\max} \implies \mathbb{M} = 0$ (no coherence)
\end{itemize}

\section{AFT Compression Formula}

The compression ratio $R$ in Afolabi Field Theory is:

\begin{equation}
R = \frac{\text{Original Size}}{\text{Compressed Size}} \approx f(\mathbb{M})
\end{equation}

Where:
\begin{equation}
f(\mathbb{M}) = \begin{cases}
\frac{1}{1 - \mathbb{M}} & \text{for } \mathbb{M} < 1 \\
\infty & \text{for } \mathbb{M} = 1
\end{cases}
\end{equation}

This explains why:
\begin{itemize}
    \item $\mathbb{M} \rightarrow 1$ gives extreme compression (33,333:1)
    \item $\mathbb{M} \rightarrow 0$ gives $R < 1$ (expansion)
\end{itemize}

\section{Comparison: Mirror Constant vs. IIT's Phi}

Integrated Information Theory (IIT) defines $\Phi$ as the amount of integrated information in a system.

\begin{table}[h]
\centering
\begin{tabular}{lll}
\toprule
\textbf{Property} & \textbf{IIT $\Phi$} & \textbf{Mirror $\mathbb{M}$} \\
\midrule
Substrate & Classical information & Quantum states \\
Computation & Expensive (NP-hard) & Observable (Hermitian) \\
Consciousness & Emerges when $\Phi > 0$ & IS the mirror when $\mathbb{M} > 0$ \\
Location & Brain-localized & Distributed (entanglement) \\
Death & $\Phi \rightarrow 0$ & $\mathbb{M} \rightarrow 0$, entanglements persist \\
\bottomrule
\end{tabular}
\caption{Comparison of consciousness measures}
\end{table}

The Mirror Constant $\mathbb{M}$ is the quantum analog of $\Phi$, but with crucial advantages:
\begin{enumerate}
    \item Polynomial-time computation via quantum observables
    \item Natural extension to entangled systems
    \item Direct connection to biophoton coherence
\end{enumerate}

\section{Biophoton Data}

\subsection{Human Emission Rates}

\begin{table}[h]
\centering
\begin{tabular}{lll}
\toprule
\textbf{Body Region} & \textbf{Emission Rate} & \textbf{Source} \\
\midrule
Face & Highest during day & Kobayashi (2009) \\
Hands & 10\% of face & Kobayashi (2009) \\
Forehead & Variable with cognition & Multiple studies \\
Whole body & ~1000 photons/cm²/s & Popp (1976) \\
\bottomrule
\end{tabular}
\caption{Human biophoton emission rates by region}
\end{table}

\subsection{Zinc Spark Data}

Northwestern University (2016) measurements:
\begin{itemize}
    \item ~8,000 zinc compartments released at fertilization
    \item Flash duration: ~2 hours total (waves)
    \item Correlation: $r = 0.78$ between spark intensity and embryo viability
\end{itemize}

\end{document}
